%%%%%%%%%%%%%%%%%%%%%%%%%%%%%%%%%%%%%%%%%%%%%%%%%%%%%%%%%%%%%%
\chapter{Conclusion}
\label{chp:conclusion}
%%%%%%%%%%%%%%%%%%%%%%%%%%%%%%%%%%%%%%%%%%%%%%%%%%%%%%%%%%%%%%
\textit{Einleitende Conclusion}

\begin{sloppypar}
    Beantworte jede Forschungsfrage direkt mit den erzielten Ergebnissen und ihrer Reichweite.
\end{sloppypar}
%%%%%%%%%%%%%%%%%%%%%%%%%%%%%%%%%%%%%%%%%%%%%%%%%%%%%%%%%%%%%%
\section{Future Work}
\label{sec:future_work}
%%%%%%%%%%%%%%%%%%%%%%%%%%%%%%%%%%%%%%%%%%%%%%%%%%%%%%%%%%%%%%

Future work should retain the reproducible discovery-and-validation protocol
while extending the mechanisms and evidence that currently limit physical and
counterfactual interpretation:

\begin{itemize}
    \item \textbf{Investigate an object-centric process representation.}
    Future work could examine whether an object-centric event log is better
    suited to container-terminal operations than the current truck-visit case
    representation. In the present model, container characteristics are
    aggregated to case level, so information such as container type, loading
    status, reefer status, or hazardous-goods status is combined across all
    containers handled during one truck visit. An object-centric approach could
    keep each container as a separate object with its own attributes while
    linking it to the truck visit, handling activity, yard block, and resource.
    This would reduce information loss and preserve container-specific information throughout the process
    and could also support multi-container and multi-block visits without
    forcing one aggregated context value onto the complete case.

    \item \textbf{Add an explicit spatial and movement representation.}
    Future work could separate physical movement from the combined
    pre-service delay by using route, position, or equipment observations.
    Rather than representing every movement as an additional process activity,
    a richer simulation formalism could include location directly in the model
    state, for example through timed or coloured Petri-net tokens and
    location-dependent transition conditions. This would allow truck or
    container movement between yard areas to consume time and respect spatial
    constraints without treating travel as a recorded handling activity.
    Movement times would then have to be estimated from spatial observations
    and the remaining access and waiting-time models refitted, so that travel
    time is not counted twice.

   \item \textbf{Represent physical assignment and queueing explicitly.}
    Future models could bind trucks and containers to physically feasible
    yard locations and resources instead of treating resource labels as fully
    interchangeable. Equipment availability and assignment should be represented
    as dynamic state, and trucks waiting for handling should enter explicit
    simulated queues rather than having these effects absorbed into a fitted
    pre-service-delay distribution. An object-centric representation could
    preserve the relations between trucks, containers, yard blocks, and
    equipment, while a corresponding simulation model would still need explicit
    movement, resource-allocation, and queueing rules
    \cite{VanDerAalst2023OCPM,Knopp2023ObjectCentricSimulation}.

    
    
\end{itemize}
